\chapter{Limits}\label{ch:limits}

This chapter says what \nwergm{} cannot do, how large a network it can handle, and which problems are known. Numbers are measured and the way they were measured is given; where nothing was measured the text says so.

\section{Network size}\label{sec:size}

\paragraph{What costs what.} The pseudolikelihood is built over \emph{all} dyads, $n(n-1)/2$ in an undirected network and $n(n-1)$ in a directed one, so that time and memory of a pseudolikelihood fit grow with the square of the number of nodes, whatever the number of ties. The Monte Carlo sampler works on the ties and their neighborhoods and its cost per proposal depends on the degrees, not on $n^2$, but every Monte Carlo fit starts from a pseudolikelihood fit (unless all coefficients are given in \opt{init()}) and so pays the cost of the dyad-level fit once, and the log-likelihood evaluation by bridge sampling needs further chains. A network is held in a sparse structure; dense $n\times n$ objects exist only for dyadic covariates (\opt{edgecov()}, \opt{hamming()}, \opt{dyadcov()}) and for \texttt{predict, matrix()}.

\paragraph{Measured run times.}\label{sec:runtimes} Run times depend on the model, the network, the settings and the load of the machine, and they changed with version 1.3.0, whose default sampling is that of \ergm{} (Section~\ref{sec:adaptive}): the sample of an iteration has a few hundred draws instead of a fixed 3000, and the fit stops with \ergm{}'s test. Table~\ref{tab:bench} compares \ergm{}~\ergmver{}, \nwergm{}~1.2.3 and \nwergm{}~1.3.0 on eight models, each with its default settings: median of five fits (seeds 1 to 5) on identical networks, the three programs run at the same time on the same machine. The measurement was made on 26 October 2026 on an Apple-silicon Mac with 18 cores that was shared with other jobs (load average 12 to 14 during the runs), with R~4.6.0 and Stata/BE~18; the times are indications of the ratio, not of absolute speed. The three undirected models were loaded as undirected networks for both programs. On these models \nwergm{}~1.3.0 was 2.4 to 7.7 times faster than \ergm{}; version 1.2.3, which recorded 3000 draws at the full interval in every iteration and drew further samples for the covariance and the log-likelihood, was slower than \ergm{} on six of the eight, about equal on one and faster on one. The speed gain comes mainly from the sample-size policy and in addition from faster compiled kernels, a frame handed to the plugin that is sized to the network, and the pseudolikelihood of large networks computed on the distinct rows of the design matrix; the estimates of the fixed-size algorithm, \opt{mcmleeffectivesize(0)}, are unchanged. The agreement of the estimates with \ergm{}'s is the subject of Chapter~\ref{ch:validation}. Speed relative to \ergm{} depends on the model, and the table does not say that every model is faster.

\begin{table}[ht]\centering\small
\caption{Median run time in seconds of five fits with default settings (26 October 2026, shared 18-core Mac, load average 12 to 14; the three programs were run at the same time). Ratio: \ergm{} divided by \nwergm{} 1.3.0. Iterations: the range of MCMLE iterations over the five fits, for \ergm{} and for \nwergm{} 1.3.0.}\label{tab:bench}
\begin{tabular}{@{}>{\raggedright\arraybackslash}p{5.3cm}rrrrl@{}}\toprule
Model & \ergm{} & 1.2.3 & 1.3.0 & ratio & iterations\\
 & & \multicolumn{2}{c}{\nwergm{}} & \ergm{}$/$1.3.0 & (R; 1.3.0)\\\midrule
30 nodes, directed: edges, mutual & 0.30 & 0.85 & 0.06 & 5.6 & 1; 1\\
100 nodes, directed: + nodematch & 0.48 & 0.94 & 0.10 & 5.0 & 1--3; 1\\
100 nodes, undirected: edges, gwesp & 0.41 & 2.31 & 0.10 & 4.2 & 1; 1\\
500 nodes, sparse undirected: + nodematch, gwesp & 3.62 & 3.74 & 0.48 & 7.7 & 3--5; 1--3\\
1000 nodes, directed: edges, mutual, nodematch & 10.46 & 6.11 & 1.86 & 5.6 & 2--3; 1--4\\
1000 nodes, sparse undirected: edges, gwesp, gwnsp & 12.30 & 22.25 & 3.52 & 3.4 & 2--6; 1--4\\
Lazega advice, 71 nodes: seven terms & 11.33 & 24.81 & 4.21 & 2.7 & 6--8; 4--12\\
Faux Mesa High, 205 nodes: gwesp(0.25) & 4.58 & 14.88 & 1.90 & 2.4 & 5--11; 5--11\\
\bottomrule\end{tabular}
\end{table}

Table~\ref{tab:scaling} gives the run time of \texttt{nwergm} on random undirected networks with mean degree 6 and three random groups, with the default settings and the native plugin, one fit per row (so the figures are indications, not benchmarks). The networks were generated by \texttt{manual/ex/ex98\_scaling.do} and \texttt{ex99\_scaling2.do}; the machine was an Apple-silicon Mac running other jobs at the same time load average about 6 during the runs of this edition, so that the times are longer than on an idle machine. The dyad-independent model has the terms \opt{edges}, \opt{nodematch(g)} and \opt{nodecov(x)}; the dyad-dependent model adds \opt{gwesp(0.5)} (Monte Carlo maximum likelihood with the default settings, with bridge sampling for the log-likelihood). The times of the pseudolikelihood column do not depend on the Markov chain; they vary between runs by tens of percent on a machine that is also used for other jobs.

\begin{table}[ht]\centering\small
\caption{Run time in seconds of \texttt{nwergm} on random networks of $n$ nodes (mean degree 6).}\label{tab:scaling}
\begin{tabular}{@{}rrrr@{}}\toprule
$n$ & dyads & pseudolikelihood (3 terms) & Monte Carlo (4 terms with \texttt{gwesp})\\\midrule
200 & 19,900 & 0.1 & 0.2 (1 iteration)\\
500 & 124,750 & 0.2 & 0.7 (1 iteration)\\
1000 & 499,500 & 0.8 & 1.6 (3 iterations)\\
2000 & 1,999,000 & 3.6 & 6.0 (3 iterations)\\
4000 & 7,998,000 & 15.7 & not run\\
5000 & 12,497,500 & not run & 39.5 (5 iterations)\\
8000 & 31,996,000 & 68.2 & not run\\
10000 & 49,995,000 & not run & 152.7 (3 iterations)\\\bottomrule
\end{tabular}
\end{table}

A random network has no clustering, so that the shared-partner statistics are small and the sampler is fast; a network with high degrees and heavy clustering is costlier per proposal and mixes more slowly, and the time of a Monte Carlo fit is dominated by the number of iterations and by the adaptive growth of the interval. These figures are \emph{not} a statement about such networks. For the Lazega network of Chapter~\ref{ch:examples} (71 nodes, 575 ties, five terms) Table~\ref{tab:timing} gives the run times of five seeds each: the default of 1.3, the fixed-size algorithm of the series 1.2 (\opt{mcmleeffectivesize(0)}, 3000 draws per iteration), the fixed-size algorithm with the settings of the versions before 1.2.0, the Mata sampler and \ergm{}.

\begin{table}[ht]\centering\small
\caption{Run time in seconds of the Lazega model (five seeds). nwergm with default settings, with fixed-size samples (as in the series 1.2) and with fixed-size samples and the settings of version 1.1 (native plugin), the Mata sampler with default settings (one run), and R's \texttt{ergm()} with its defaults, on the same Mac with other jobs running (load average about 6).}\label{tab:timing}
\begin{tabular}{@{}lrrrr@{}}\toprule
setting & runs & median & range & iterations\\\midrule
nwergm, default settings & 5 & 0.7 & 0.7--0.8 & 1--1\\
nwergm, \texttt{mcmcinterval(1000) mcmcburnin(20000)} & 5 & 6.9 & 6.9--7.0 & 1--1\\
nwergm, Mata sampler (\texttt{nonative}, default settings) & 1 & 34.0 & 34.0--34.0 & 1--1\\
R \texttt{ergm()}, default settings & 5 & 11.2 & 5.7--16.6 & 3--8\\
\bottomrule\end{tabular}

\end{table}

Run times of \nwergm{} and \ergm{} depend on the model, the network, the settings and the load of the machine; two programs with different stopping rules and chain lengths are not compared by one example, and the table of eight models above is the better guide. On the Lazega model of this table \nwergm{} 1.3.0 took a median of 1.2 seconds (2 to 7 iterations) and \ergm{} 10.8 seconds (3 to 8 iterations), measured one after the other in the same session, so that the load of the machine was not the same for the two; the same model fitted with fixed 3000 draws, as in the series 1.2, took 2.3 seconds in one iteration for every seed. The Mata sampler (\opt{nonative}) is the slow fallback: 262 seconds for one fit. The Monte Carlo error of the two programs is similar (Chapter~\ref{ch:validation}).

\paragraph{Larger networks.} Not measured: networks above 10,000 nodes, networks with mean degree above 6, models with many terms or with several dense dyadic covariates, memory use as a function of $n$, and fits of a directed network of the sizes in the table.

\paragraph{Stata limits.} Stata/BE has a maximum matrix dimension of 800 (\texttt{c(max\_matdim)}) and at most 2048 variables. A network with more than 800 nodes therefore cannot be declared from a Stata matrix with \texttt{nwset, mat()}; build it from an edge list (\texttt{nwfromedge}), as the networks of Table~\ref{tab:scaling} were. The same holds for \texttt{predict, matrix()}; \texttt{predict, frame()} has one row per dyad and is limited by memory. The matrix \texttt{e(mcmcsample)} has one row per draw; filling it from Mata is the slowest step for a very large \opt{mcmcsamplesize()} (more than 30 seconds at 100,000 rows in a test of the development of the package), and \opt{nomcmcsample} skips it. Other editions of Stata have larger limits; they were not tested.

\paragraph{Capacity of the plugin.} The plugin has fixed capacities: 256 coefficients (columns of the model), 64 distinct node attributes and 8 distinct dense dyadic covariate matrices in one model. A model that exceeds them is not refused; it runs in the Mata sampler as a whole, never partly in C and partly in Mata, and is much slower (and \opt{structural()}, \opt{operators()} and the constraints, which exist only in the plugin, are refused).

\section{The plugin and the platforms}\label{sec:platforms}

The native plugin is built for macOS (universal binary, compiled on the development machine), Linux and Windows (compiled by continuous integration). All the output of this manual comes from the macOS build with Stata/BE~18; the Linux and Windows plugins were not run for this manual, and the other editions of Stata (SE, MP) were not tested. The plugin reports its version, and a plugin that is older than the \texttt{.ado} files is detected. A platform without a plugin runs the Mata sampler, which is much slower for the models that have one (Table~\ref{tab:timing}) and does not exist for \opt{structural()}, \opt{operators()}, \opt{constraints()}, \opt{obsconstraints()} and \opt{mcmcmaxedges()}.

Two practical traps. A plugin in Stata's PLUS directory from an older installation shadows the plugin of a repository on the ado-path, so that a freshly built plugin can silently not be used; test it through a scratch ado-path. On macOS, overwriting a plugin that has already been loaded kills Stata without a message (the code signature is cached by the system); remove the old file first, then copy.

\section{Dialog boxes}

The three dialogs are generated from one specification and checked automatically: a script builds the command that the dialog would build for every term and option and runs it in Stata. They were not tested by hand on Windows or Linux, and the screen layout was not tested on other platforms than macOS. A dialog builds a single command line; model specifications that are easier to write by hand (long operator formulas, \opt{For()} loops) are better typed in a do-file.

\section{Known issues}\label{sec:issues}

\begin{itemize}
  \item \textbf{Valued networks are fitted as binary graphs.} A network flagged valued is fitted to the graph in which a tie is present when its value is not zero, with one note; the values are not an outcome of the model and there is no valued ERGM (Section~\ref{sec:data}, Section~\ref{sec:types}). Temporal networks are refused.
  \item \textbf{Changed defaults in versions 1.2.0 and 1.3.0.} Version 1.3.0 changed the default sample size and the convergence test to those of \ergm{} (adaptive sample, target effective size 64; Section~\ref{sec:adaptive}): a fit with default settings draws different, statistically equivalent samples than in the series 1.2 (the estimates differ by Monte Carlo error, the number of draws is \texttt{e(mcmc\_samplesize)} and no longer 3000), is several times faster (Table~\ref{tab:bench}), and its log-likelihood has a somewhat larger Monte Carlo standard error \texttt{e(ll\_se)} because the bridges use the smaller sample of the fit. \opt{mcmleeffectivesize(0)} gives the fixed-size algorithm of 1.2. Results printed in earlier versions of this manual, in the help files or in the book for fits with default settings differ in the last digits or more. Version 1.2.0 had changed the defaults before: The default burn-in and interval of the Markov chain are those of \ergm{} (16384 and 1024; before: 3000 and 50), so that a fit with default settings gives other numbers than in version 1.1.x (the same seed gives another chain; fits that give both options are unchanged) and mixes better (Section~\ref{sec:mcmcopts}). The old settings are \opt{mcmleeffectivesize(0) mcmcburnin(3000) mcmcinterval(50)}.
  \item \textbf{Degenerate and poorly mixing models.} A model that the data do not support makes the Markov chain degenerate. The fit then stops with the error of \ergm{} (Section~\ref{sec:kapferer}) or, worse, converges although the chain mixes poorly (Section~\ref{sec:lazega}). Once the chain does not mix, the adaptive sample is extended and thinned in every iteration (at most 16 runs), a failed convergence test raises the target and the base size (up to 32 times their starting values), and every iteration costs a multiple of the first; the default 60 iterations then take many minutes. Lower \opt{mcmleiterations()} for a model that you suspect to be degenerate, and always look at \texttt{estat mcmcdiag}.
  \item \textbf{The Mata sampler can be extremely slow.} \nwergm{} prints one warning line before a Monte Carlo fit, \texttt{estat gof} or \texttt{nwergm simulate} that runs on the Mata sampler instead of the plugin, with the reason (the plugin file is missing for the platform, \opt{nonative}, or a model outside the capacity of the plugin); no warning is printed when the plugin is used. For the Lazega model with the default settings it took 262 seconds against 1.2 seconds with the plugin (Table~\ref{tab:timing}); a fit of \opt{bgwdsp1(0.5)} on the 42-node two-mode network of Section~\ref{sec:twomode} with \opt{nonative} and \opt{mcmcsamplesize(500)} had not finished after 29 minutes of processor time and was stopped (a test during the development of the package, not repeated for this release). \opt{nonative} is for testing and comparison.
  \item \textbf{Dyad-dependent constraints are slow.} A fit of \opt{edges mutual nodematch(prac)} to the 71-node Lazega network with \opt{constraints(bd(maxout=30))} took 3.8 seconds with the default settings, against 0.2 seconds without the constraint (the bound does not bind: the largest out-degree is 22; both fits needed four iterations; \texttt{manual/ex/ex9d\_bdtime.do}; with the series 1.2 the same fit took 191 seconds). Dyad-dependent constraints change the proposal and the starting values, and the time of such a fit is not predictable from that of the unconstrained model.
  \item \textbf{Standard errors of a pseudolikelihood fit of a dyad-dependent model} are not valid (Section~\ref{sec:mple}); \opt{mplecovariance(godambe)} or \opt{mplecovariance(bootstrap)} give covariances for the pseudolikelihood estimate.
  \item \textbf{No likelihood for a curved Monte Carlo fit.} \texttt{e(ll)}, AIC and BIC are missing after a fit with a free decay by Monte Carlo maximum likelihood (the pseudolikelihood fit has them), and \texttt{estat ic} refuses.
  \item \textbf{Missing dyads and constraints.} Missing dyads cannot be combined with \opt{freedyads()}, \opt{blockdiag()} or \opt{fixdensity}.
  \item \textbf{Collinear terms.} A set of terms that is exactly collinear (for example all levels of \opt{nodefactor()} with \opt{base(0)} together with \opt{edges}) is reported with a coefficient of 0 and no standard error; a statistic that never varies in the sample produces the warning that the model may be nonidentifiable.
  \item \textbf{Operators.} Operators cannot be combined with constraints, and \texttt{S()} and \texttt{Project()} do not accept dyadic covariates (Section~\ref{sec:operators}).
\end{itemize}

\section{Messages and what to do}\label{sec:messages}

The messages that the examples of this manual produced, and the usual remedy.
\footnotesize
\begin{longtable}{@{}>{\raggedright\arraybackslash}p{6.4cm}>{\raggedright\arraybackslash}p{9.3cm}@{}}
\toprule
Message (return code) & Meaning and remedy\\
\midrule\endhead
\texttt{option edges is required} (198) & Every model has the term \opt{edges}; \opt{structural(density)} or \opt{structural(meandeg)} can stand in for it.\\
\texttt{option base() not allowed} (198) & A suboption was typed outside the parentheses of its term. Write \opt{nodefactor(x, base(0))}, not \opt{nodefactor(x) base(0)}.\\
\texttt{In term nodecov: attribute wealth has missing data, which is not currently supported.} (416) & A node attribute that a term uses has a missing value, as in \ergm{}. Drop or impute it (Section~\ref{sec:missing}).\\
\texttt{Unconstrained MCMC sampling did not mix at all. Optimization cannot continue.} (498) & The statistics of the sample do not vary: the model is degenerate or the starting values are far from any fit. Simplify the model, replace a count term by its geometrically weighted version, or give starting values with \opt{init()} (Section~\ref{sec:kapferer}).\\
\texttt{Warning: MCMLE iter 1 - linear dependence detected among the model statistics} & Some statistics are collinear in the sample, for example a term that is constant under a constraint or is a function of other terms. The model may be nonidentifiable; remove the term.\\
\texttt{Note: the specified model's sample space constraint holds statistic(s) edges constant. They will be ignored.} & As in \ergm{}: a statistic that the constraint fixes is shown as omitted (Section~\ref{sec:constraints}).\\
\texttt{Note: method(mple) requested for a dyad-dependent model} & The pseudolikelihood is not the likelihood and its standard errors are not valid (Section~\ref{sec:mple}).\\
\texttt{The MPLE does not exist!} (warning) & The data are separated for this model. Use \opt{mpletype(penalized)} or drop the term (Section~\ref{sec:mple}).\\
\texttt{Note: network \ldots{} has 3 missing dyad(s)} & Missing dyads are excluded from the likelihood and handled as in \ergm{} (Section~\ref{sec:missing}).\\
\texttt{estat mcmcdiag is only available after method(mcmle)} (498) & There is no Markov chain after a pseudolikelihood fit; also not available after \opt{nomcmcsample}.\\
\texttt{\ldots{} need the native plugin} (198) & The option exists only in the plugin, and no plugin is available for the model or the platform, or \opt{nonative} was given (Section~\ref{sec:platforms}).\\
\texttt{The installed core nwcommands and nwergm were built for different versions} (3021 or 3260) & The Mata libraries of the core and of \nwergm{} do not match. Reinstall the models with \texttt{nwinstall, models} or all packages with \texttt{nwinstall, everything} (Section~\ref{sec:platforms}).\\
\bottomrule
\end{longtable}
\normalsize
